\section{Dimension Reduction}

\begin{frame}{Tổng quan về giảm chiều dữ liệu}
	\begin{block}{Định nghĩa}
		Giảm chiều dữ liệu (dimensionality reduction) là quá trình biến đổi dữ liệu từ không gian đặc trưng ban đầu $P$-chiều
		sang không gian biểu diễn thấp chiều hơn $K$-chiều, với $K \ll P$.
	\end{block}

	\vspace{1em}
	% Mục tiêu của giảm chiều là giúp trực quan hóa dữ liệu tốt hơn, giảm nhiễu do các chiều kém hữu ích, và trích xuất biểu diễn mới phù hợp hơn cho các mô hình downstream.

	Với dữ liệu $\mathbf{X} \in \mathbb{R}^{N \times P}$, ta tìm ánh xạ
	\[
		f: \mathbb{R}^{P} \rightarrow \mathbb{R}^{K}, \qquad z_i = f(x_i)
	\]
	sao cho vẫn giữ được cấu trúc quan trọng của dữ liệu.
	\hfill{\scriptsize Theo \cite{duchesnay2021statsml}}
\end{frame}
\begin{frame}{Mục tiêu của giảm chiều dữ liệu}
    Giảm chiều dữ liệu không chỉ để nén dữ liệu, mà còn nhằm giữ lại cấu trúc quan trọng
    phục vụ mô hình hóa và phân tích.

    \vspace{0.6em}
    \begin{enumerate}
        \item \textbf{Trực quan hóa}: Chiếu dữ liệu về 2D/3D để quan sát phân bố, xu hướng và cụm dữ liệu.
        \item \textbf{Giảm nhiễu}: Loại bỏ chiều kém hữu ích, tăng độ ổn định và độ chính xác của mô hình.
        \item \textbf{Trích xuất đặc trưng}: Tạo biểu diễn mới chất lượng hơn cho các mô hình downstream.
    \end{enumerate}

    \vspace{0.4em}
    \begin{block}{Kết luận}
        Các phương pháp như PCA bản chất là kỹ thuật trích xuất đặc trưng, chứ không chỉ là công cụ nén dữ liệu.
    \end{block}

    \hfill{\scriptsize Theo \cite{duchesnay2021statsml}}
\end{frame}
\begin{frame}{Phân loại hai nhóm phương pháp giảm chiều}
    \begin{columns}[T,onlytextwidth]
        \begin{column}{0.48\textwidth}
            \textbf{Phương pháp tuyến tính (Linear)}
            \vspace{0.4em}
            \begin{itemize}
                \item Dựa trên biến đổi tuyến tính: PCA, SVD.
                \item Khai thác ma trận hiệp phương sai giữa đặc trưng.
                \item Phù hợp khi dữ liệu có cấu trúc gần tuyến tính.
                \item Thường dùng làm bước tiền xử lý cho mô hình dự đoán.
            \end{itemize}
        \end{column}

        \begin{column}{0.48\textwidth}
            \textbf{Phương pháp phi tuyến (Non-linear)}
            \vspace{0.4em}
            \begin{itemize}
                \item Dựa trên manifold learning: Isomap, t-SNE, MDS.
                \item Tìm embedding thấp chiều cho dữ liệu cong/phức tạp.
                \item Bảo toàn cấu trúc lân cận/cấu trúc hình học phi tuyến.
                \item Thường dùng cho khám phá dữ liệu và trực quan hóa.
            \end{itemize}
        \end{column}
    \end{columns}

    \vspace{0.3em}
    \begin{alertblock}{Ghi chú}
        Hai nhóm phương pháp mang tính bổ trợ, không thay thế nhau hoàn toàn.
    \end{alertblock}
\end{frame}

\begin{frame}{Vai trò của giảm chiều trong pipeline Machine Learning}
    \centering
    \begin{tikzpicture}[
        >=stealth,
        pipestep/.style={rectangle, rounded corners, draw=black!70, fill=blue!8, align=center, minimum width=2.4cm, minimum height=0.8cm, font=\scriptsize},
        methodbox/.style={rectangle, rounded corners, draw=black!70, fill=green!10, align=center, minimum width=2.7cm, minimum height=0.9cm, font=\scriptsize},
        notebox/.style={rectangle, rounded corners, draw=black!60, fill=orange!10, align=center, minimum width=3.2cm, minimum height=0.8cm, font=\scriptsize}
    ]
        % Pipeline steps
        \node[pipestep] (s1) at (0,0) {Thu thập\\dữ liệu};
        \node[pipestep] (s2) at (2.7,0) {Tiền xử lý};
        \node[pipestep] (s3) at (5.4,0) {Giảm chiều};
        \node[pipestep] (s4) at (8.1,0) {Huấn luyện};
        \node[pipestep] (s5) at (10.8,0) {Đánh giá};

        \draw[->, thick] (s1) -- (s2);
        \draw[->, thick] (s2) -- (s3);
        \draw[->, thick] (s3) -- (s4);
        \draw[->, thick] (s4) -- (s5);

        % Methods
        \node[methodbox] (lin) at (5.4,1.7) {PCA / SVD\\(tuyến tính)};
        \node[methodbox] (nonlin) at (5.4,-1.7) {Isomap, t-SNE, MDS\\(phi tuyến)};
        \node[notebox] (viz) at (9.2,-1.7) {Khám phá dữ liệu\\và trực quan hóa};

        % Method-to-step links
        \draw[->, thick, green!50!black] (lin) -- (s3);
        \draw[->, thick, green!50!black] (lin.east) to[bend left=12] (s4.north);

        \draw[->, thick, green!50!black] (nonlin) -- (s3);
        \draw[->, thick, green!50!black] (nonlin.east) -- (viz.west);
    \end{tikzpicture}

\end{frame}